\setcounter{paragraphcounter}{1}

\section{Satzung}

\noindent\textbf{Anmerkung:} Zur besseren Lesbarkeit werden alle Personentitel in der maskulinen Form aufgeführt, gelten aber auch für die feminine und die diverse Form.

\subsection{§ \theparagraphcounter\ Name, Sitz und Geschäftsjahr}

\begin{enumerate}
    \item Der Verein führt den Namen "Narrenbund Ostfildern e.V."
    
    \item Der Verein ist unter der Registernummer VR 211 613 im Vereinsregister Stuttgart eingetragen und hat seinen Sitz in Ostfildern-Nellingen.
    
    \item Der Verwaltungssitz und Vereinsanschrift ist die Postanschrift des jeweiligen 1. Vorstands; Änderungen führen nicht zu einer Satzungsänderung.
    
    \item Das Geschäftsjahr des Vereins ist das Kalenderjahr.
    
    \item Der Verein kann weiteren Abteilungen einen eigenen Gruppennamen geben, der aber vom Vereinsausschuss genehmigt wird.
\end{enumerate}

\refstepcounter{paragraphcounter}
\subsection{§ \theparagraphcounter\ Zweck des Vereins}

\begin{enumerate}
    \item Der Zweck des Vereins ist die Pflege und Förderung des Fasnetbrauchtums im Sinne des § 52 Abs. 2 Nr 22 Abgabenordnung (AO). Es ist zugrunde gelegt, dass in Ostfildern die schwäbisch-alemannische Fastnacht bodenständig ist.
    \\
    Der Vereinszweck wird vor allem durch folgende Aktivitäten verwirklicht:
    \begin{itemize}
        \item Unterhalt von Maskengruppen.
        \item Teilnahme an bzw. Durchführung von Fastnachtsveranstaltungen.
        \item Unterstützung und Zusammenarbeit mit anderen Vereinen im Bereich der schwäbisch-alemannische Fastnacht.
        \item Die ganzjährige Förderung von Tanz und Spiel, insbesondere der Jugendlichen.
    \end{itemize}
    
    \item Zur schwäbisch-alemannische Fastnacht gehören die Maskengruppen. Über weitere Zulassungen von Maskengruppen und sonstiger Gruppen entscheidet die Mitgliederversammlung.

    \item Bestrebungen parteipolitischer, konfessioneller und rassistischer Art sind im Verein ausgeschlossen.

    \item Der Verein ist selbstlos tätig und verfolgt nicht in erster Linie einen eigenwirtschaftlichen Zweck.
        
    \item Der Verein verfolgt ausschließlich und unmittelbar gemeinnützige Zwecke im Sinne des §§ 51-68 Abgabenordnung (AO). Etwaige Überschüsse dürfen nur für satzungsgemäße Zwecke verwendet werden. Ansammlung von Vermögen zu anderen Zwecken ist nicht zulässig. Die Mitglieder erhalten keine Gewinnanteile und in ihrer Eigenschaft als Mitglieder auch keine sonstigen Zuwendungen aus Mitteln des Vereins. Bei ihrem Ausscheiden oder bei Auflösung des Vereins erhalten sie weder einbezahlte Beiträge zurück noch habe sie einen Anspruch auf Vereinsvermögen. Es darf keine Person durch Ausgaben, die dem Zweck des Vereins fremd sind, oder durch unverhältnismäßig hohe Vergütungen begünstigt werden. Die Organe des Vereins arbeiten ehrenamtilich. Für entstandene Ausgaben wird Ersatz geleistet.
\end{enumerate}

\refstepcounter{paragraphcounter}
\subsection{§ \theparagraphcounter\ Mitgliedschaft}
\begin{enumerate}
    \item Der Verein gliedert sich in:
    \begin{enumerate}
        \item ordentliche Mitglieder (aktive Mitglieder)
        \item außerordentliche Mitglieder (passive Mitglieder)
        \item Ehrenmitglieder
    \end{enumerate}
\end{enumerate}

\refstepcounter{paragraphcounter}
\subsection{§ \theparagraphcounter\ Erwerb und Erlöschen der Mitgliedschaft}
\begin{enumerate}
    \item Mitglied des Narrenbund Ostfildern e.V. kann jede natürliche und juristische Person werden, die die Ziele des Vereins unterstützen und die Satzung sowie bestehende Vereinsordnung akzeptieren.

    \item Juristische Personen können nur als außerordentliche Mitglieder aufgenommen werden.

    \item Zu Ehrenmitgliedern können Personen ernannt werden, die sich besondere Verdienste erworben haben. Die Ernennung erfolgt durch Beschluss der Mitgliederversammlung mit 2/3 Mehrheit der gültig abgegebenen Stimmen.
    
    \item Der schriftliche Aufnahmeantrag ist an den geschäftsführenden Vorstand des Vereins zu richten, wobei Minderjährige eine Zustimmungserklärung ihres gesetzlichen Vertreters bedürfen.

    \item Über die Aufnahme als Probeläufer entscheidet der Vereinsausschuss.

    \item Die Mitgliedschaft wird durch Eintritt in den Verein erlangt. Über die Aufnahme neuer Mitglieder entscheidet nach Abschluss des Probejahres der Vereinsausschuss.

    \item Mitglieder, die im Laufe ihrer Mitgliedschaft ihren Mitgliedsstatus von aktiv auf passiv setzen und nach unbestimmter Zeit wieder in den aktiven Status zurückwechseln, können dies ohne Einschränkungen tun.

    \item Bei passiven Mitgliedern, die noch kein Probejahr absolviert haben und in den aktiven Mitgliederstatus wechseln wollen, wird wie bei einer Neumitgliedschaft verfahren.

    \item Eine eventuelle Ablehnung eines Aufnahmeantrages bedarf keiner Begründung, es besteht auch kein Anspruch des Antragstellers auf Begründung der Ablehnung.

    \item Die Mitgliedschaft erlischt durch Austritt, Ausschluss, Tod oder Auflösung des Vereins.

    \item Der Austritt eines Mitgliedes kann bis zu drei Monate zum Ende eines Kalenderjahres durch schriftliche Mitteilung an den Vorstand des Vereins erfolgen. Für die Austrittserklärung Minderjähriger gelten die für den Aufnahmeantrag bestimmten Regelungen entsprechend. Die finanziellen Verpflichtungen für das laufende Kalenderjahr werden durch das Ausscheiden nicht berührt.

    \item Mitglieder können ggf. durch einfachen Mehrheitsbeschluss des Vereinsausschusses ausgeschlossen werden bei:
    \begin{enumerate}
        \item groben und wiederholten Verstößen gegen die Satzung oder gegen satzungsgemäße Beschlüsse.

        \item unehrenhaftem Verhalten und Verlust der bürgerlichen Ehrenrechte.

        \item nicht nachkommen der mitgliedschaftlichen Verpflichtungen.

        \item nichtleistung fälliger Beitragszahlungen trotz Mahnung mit einer Frist von vier Wochen.
    \end{enumerate}
    Vor Beschlussfassung ist dem betroffenen Mitglied Gelegenheit zur Äußerung zu geben. Der Ausschließungsbeschluss muss unter Angabe der Gründe, die zum Ausschluss führten, dem Mitglied schriftlich mitgeteilt werden. Gegen den Ausschluss kann das betroffene Mitglied innerhalb eines Monats schriftlich Widerspruch einlegen. Über den Widerspruch entscheidet die Mitgliederversammlung mit einfacher Mehrheit. Der Ausschluss wird bis zur Entscheidung der Mitgliederversammlung ausgesetzt.

    \item Mit Beendigung der Mitgliedschaft erlöschen alle Ansprüche aus dem Mitgliedsverhältnis, unbeschadet des Anspruches des Vereins auf bestehende Forderungen.
\end{enumerate}

\refstepcounter{paragraphcounter}
\subsection{§ \theparagraphcounter\ Mitgliedsbeitrag}
\begin{enumerate}
    \item Von den Mitgliedern werden Beiträge erhoben. Der jährliche Mitgliedsbeitrag wird auf Vorschlag des Vereinsausschusses von der Mitgliederversammlung beschlossen.

    \item Der Mitgliedsbeitrag ist beim Eintritt in den Verein und jährlich im 2. Quartal des Geschäftsjahrs fällig.

    \item Ehrenmitglieder sind von der Beitragspflicht befreit.
\end{enumerate}

\refstepcounter{paragraphcounter}
\subsection{§ \theparagraphcounter\ Rechte und Pflichten der Mitglieder}
\begin{enumerate}
    \item Für die Mitglieder ist diese Satzung, die Geschäftsordnung sowie die Beschlüsse der Vereinsorgane verbindlich.

    \item Mit dem Beitritt als ordentliches Mitglied des Narrenbund Ostfildern e.V., erkennt man die damit einhergehenden Pflichten an. Dazu gehört es der Förderungspflicht nachzukommen, welche sich u.a. auch in der Beteiligung bei Arbeitsdiensten ausdrückt. Art und Umfang der Arbeitsdienste können in einer Ordnung des Vereins geregelt werden. Dabei sind persönliche Einschränkungen der Mitglieder angemessen zu berücksichtigen.

    \item Jedes Mitglied ist zur Zahlung des Mitgliedsbeitrages sowie zu sonstigen beschlossenen oder durch die Satzung/sonstige Ordnung festgelegte Abgaben/Umlagen verpflichtet. 

    \item Die Mitglieder sind verpflichtet, die Vereinsinteressen zu fördern und alles zu unterlassen, was dem Ansehen und dem Zweck des Vereins entgegensteht.

    \item Ordentliche, außerordentliche und Ehrenmitglieder, die das 18. Lebensjahr vollendet haben, sind berechtigt, an der Willensbildung im Verein durch Ausübung des Antrags-, Diskussions- und Stimmrechts in der Mitgliederversammlung teilzunehmen. Jedes stimmberechtigte Mitglied hat eine Stimme.

    \item Mitglieder, die das 18. Lebensjahr noch nicht vollendet haben, besitzen kein Stimm- und kein aktives oder passives Wahlrecht.
\end{enumerate}

\refstepcounter{paragraphcounter}
\subsection{§ \theparagraphcounter\ Organe des Vereins}
\begin{enumerate}
    \item Die Organe des Vereins sind:
    \begin{enumerate}
        \item Der Vorstand
        \item Der Vereinsausschuss
        \item Die Mitgliederversammlung
    \end{enumerate}
\end{enumerate}

\refstepcounter{paragraphcounter}
\subsection{§ \theparagraphcounter\ Der Vorstand}
\begin{enumerate}
    \item Den Vorstand bilden:
    \begin{enumerate}
        \item der 1. Vorstand (Oberzunftmeister)
        \item der 1. stellv. Vorstand (Zunftmeister)
        \item der 2. stellv. Vorstand bzw. Kassier (Schatzmeister)
    \end{enumerate}
    
    \item Sie sind die gesetzlichen Vertreter im Sinne des § 26 BGB.

    \item Jeweils zwei dieser Vorstandsmitgliewder vertreten den Verein gemeinsam.

    \item Der Vorstand ist berechtigt, Ausgaben bis zu 1.000 EUR pro Einzelfall eigenverantwortlich zu tätigen. Beträge über 1.000 EUR bedürfen der vorherigen Zustimmung des Vereinsausschusses. Über alle Ausgaben ist der Ausschuss unverzüglich, spätestens jedoch innerhalb von einer Woche, zu informieren.

    \item Der Vorstand ist  im Haftungsfall vom Haftungsanspruch freigestellt (gem. § 31 Abs. 2 BGB). Dieses gilt jedoch nicht bei grober Fahrlässigkeit oder Vorsatz.

    \item Der Vorstand ist für alle Angelegenheiten des Vereins zuständig, soweit sie nicht durch die Satzung einem anderen Organ zugewiesen sind. Zu seinen Aufgaben zählen insbesondere die:
    \begin{itemize}
        \item Vorbereitung und Einberufung der Mitgliederversammlung und Aufstellung der Tagesordnung.
        
        \item Leitung bei allen Versammlungen.

        \item Ausführung sämtlicher Beschlüsse des Narrenbund Ostfildern e.V.

        \item Aufstellung eines etwaigen Haushaltsplans und Erstellung eines Jahresberichts.
    \end{itemize}

    \item Die von der Mitgliederversammlung zu wählenden Mitglieder des Vorstandes werden auf die Dauer von drei Jahren gewählt. Sie bleiben bis zur Bestellung des neuen Vorstandes im Amt.

    \item Wählbar sind nur voll aufgenommene ordentliche Mitglieder aus den jeweiligen Gruppen, keine Probeläufer. Die doppelte Belegung von Ämtern durch eine Person innerhalb des Vorstandes ist nicht zulässig. 

    \item Scheidet ein Mitglied des Vorstandes während der Amtsperiode aus, so wird eine außerordentliche Mitgliederversammlung zur Wahl des entsprechenden Amtes vom Vorstand einberufen.

    \item Die Beschlüsse sind schriftlich zu protokollieren und von dem Vorstandsvorsitzenden zu unterzeichnen.

    \item Die Zuständigkeiten innerhalb des Vorstands sind in der Geschäftsordnung definiert.
\end{enumerate}

\refstepcounter{paragraphcounter}
\subsection{§ \theparagraphcounter\ Der Vereinsausschuss}
\begin{enumerate}
    \item Der Vereinsausschuss besteht aus:
    \begin{enumerate}
        \item dem Vorstand
        \item dem Schriftführer
        \item dem Festwart
        \item dem Häswart
        \item dem Jugendleiter
    \end{enumerate}
    
    \item Der Vereinsausschuss ist beschlussfähig, wenn mindestens 50 \% der Ausschussmitglieder anwesend sind.

    \item Die Beschlüsse werden mit einfacher Mehrheit gefasst. Jedes Mitglied des Vereinsausschusses hat eine Stimme; bei Stimmengleichheit entscheidet die Stimme des 1. Vorstands bzw. bei Abwesenheit die des jeweiligen Vertreters.

    \item Der Vereinsausschuss tritt auf Einladung des 1. Vorstands oder bei dessen Verhinderung des 1. stellv. Vorstands zusammen.

    \item Die von der Mitgliederversammlung zu wählenden Mitglieder des Vereinsausschusses werden mit der Ausnahme der Vorstände auf die Dauer von zwei Jahren gewählt. Sie bleiben bis zur Bestellung des neuen Vereinsausschusses im Amt.

    \item Wählbar sind nur voll aufgenommene ordentliche Mitglieder aus den jeweiligen Gruppen, keine Probeläufer.

    \item Der Jugendleiter wird von den Mitgliedern der Jugend gewählt. Die Wahl bedarf in jedem Fall der Bestätigung durch die Mitgliederversammlung.

    \item Scheidet während der Amtszeit ein Mitglied des Vereinsausschusses aus, so kann der Vereinsausschuss einen kommissarischen, stimmberechtigten Vertreter bis zum nächsten turnusgemäßen Wahltermin bestimmen. Diese Regelung gilt nicht für die Mitglieder des Vorstandes nach § 26 BGB.

    \item Die Beschlüsse sind schriftlich zu protokollieren und von dem Vorstandsvorsitzenden zu unterzeichnen.

    \item Die Zuständigkeiten innerhalb des Vereinsausschusses sind in der Geschäftsordnung definiert. 
\end{enumerate}

\refstepcounter{paragraphcounter}
\subsection{§ \theparagraphcounter\ Die Mitgliederversammlung}
\begin{enumerate}
    \item Die Mitgliederversammlung stellt die Richtlinien für die Arbeit des Vereins auf und entscheidet Fragen von grundsätzlicher Bedeutung. Sie wird vom 1. Vorstand, bei dessen Verhinderung von einem anderen Vorstandsmitglied geleitet. Zu den Aufgaben der Mitgliederversammlung gehören insbesondere:
    \begin{itemize}
        \item Wahl und Abberufung des Vorstandes.
        \item Wahl und Abberufung des Ausschusses.
        \item Wahl von zwei Kassenprüfern.
        \item Entgegennahme des Jahresberichts der Vorstände.
        \item Beschlussfassung über die Entlastung des Vorstandes und des Schatzmeisters.
        \item Genehmigung des Kassenberichtes und der Vermögensaufstellung des Schatzmeisters.
        \item Festsetzung der Höhe des Mitgliedsbeitrags.
        \item Beschlussfassung über die Ernennung von besonders verdienstvollen Personen zu Ehrenmitgliedern.
        \item Beschlussfassung über Aufnahme neuer Mitglieder.
        \item Beschlussfassung über Änderung der Satzung und über die Auflösung des Vereins.
    \end{itemize}

    \item Die Einberufung der Mitgliederversammlung erfolgt mindestens drei Wochen vorher per E-Mail an die Mitglieder. Zusätzlich kann die Einberufung im Amtsblatt der Stadt Ostfildern veröffentlicht und auf der Internetseite des Vereins eingestellt werden. Maßgeblich für die Wahrung der Frist ist ausschließlich die Versendung der E-Mail.

    \item Eine ordentlich einberufene Mitgliederversammlung ist grundsätzlich beschlussfähig.

    \item Die Mitgliederversammlung soll im ersten Halbjahr des neuen Geschäftsjahres abgehalten werden. Der Zeitpunkt und die Tagesordnung werden durch den Vorstand festgelegt.

    \item Sofern Wahlen durchzuführen sind, müssen diese in der Tagesordnung der Mitgliederversammlung ausdrücklich angekündigt werden.

    \item Ein Mitglied kann bis spätestens eine Woche vor dem angesetzten Termin schriftlich die Ergänzung der Tagesordnung beantragen. Die ergänzenten Tagesordnungspunkte sind zu Beginn der Mitgliederversammlung bekanntzugeben. 

    \item Über Anträge auf Ergänzung der Tagesordnung, die erst in der Hauptversammlung gestellt werden, beschließt die Hauptversammlung. Zur Annahme des Antrags ist eine Mehrheit von zwei Dritteln der abgegebenen gültigen Stimmen erforderlich.

    \item Über die Beschlüsse der Hauptversammlung ist ein Protokoll aufzunehmen, das vom jeweiligen Versammlungsleiter und dem Protokollführer zu unterzeichnen ist. 

    \item Der Protokollführer wird vom Versammlungsleiter bestimmt, in der Regel der gewählte Schriftführer.

    \item Beschlüsse der Mitgliederversammlung werden mit einfacher Mehrheit gefasst.

    \item Zur Ausübung des Stimmrechts kann ein anderes Mitglied schriftlich bevollmächtigt werden. Ein Mitglied darf jedoch nicht mehr als drei Stimmen vertreten.

    \item Zur Änderung der Satzung ist eine Mehrheit von zwei Dritteln der abgegebenen gültigen Stimmen erforderlich. Die Änderung der Satzung ist im Vereinsregister bei dem zuständigen Amtsgericht in Stuttgart einzutragen.

    \item Eine Änderung des Zwecks des Vereins kann nur mit Zustimmung aller Mitglieder beschlossen werden. Die schriftliche Zustimmung der in der Hauptversammlung nicht erschienenen Mitglieder kann nur innerhalb eines Monats gegenüber dem Vorstand erklärt werden.

    \item Jedes Mitglied kann Wahlvorschläge einreichen. Die zur Wahl vorgeschlagenen Mitglieder sind zu befragen, ob sie das Amt im Falle der Wahl annehmen. Von nicht anwesenden Kandidaten muss darüber eine schriftliche Erklärung vorliegen.

    \item Hat im ersten Wahlgang kein Kandidat die einfache Mehrheit der abgegebenen gültigen Stimmen erreicht, findet eine Stichwahl zwischen den Kandidaten statt, welche die beiden höchsten Stimmenzahlen erreicht haben. Gewählt ist dann derjenige, der die meisten gültigen Stimmen auf sich vereint. Eine Stimmengleichheit gilt als Ablehnung.

    \item Die Wahlen bzw. Abstimmungen müssen schriftlich bzw. geheim durchgeführt werden, wenn ein anwesendes, stimmberechtigtes Mitglied dies beantragt.
\end{enumerate}

\refstepcounter{paragraphcounter}
\subsection{§ \theparagraphcounter\ Außerordentliche Mitgliederversammlung}
\begin{enumerate}
    \item Eine außerordentliche Mitgliederversammlung findet statt, wenn mindestens ein viertel aller Mitglieder sie unter Angabe von Gründen schriftlich verlangen.

    \item Der Vorstand kann jederzeit eine außerordentliche Mitgliederversammlung einberufen, wenn das Interesse des Vereins es erfordert.

    \item Für die außerordentliche Mitgliederversammlung gilt die Satzung § 10.
\end{enumerate}

\refstepcounter{paragraphcounter}
\subsection{§ \theparagraphcounter\ Kassenprüfung}
\begin{enumerate}
    \item Die Mitgliederversammlung wählt zwei Kassenprüfer für die Dauer von zwei Jahren. Wiederwahl ist zulässig.

    \item Kassenprüfer dürfen weder dem Vorstand noch dem Vereinsausschuss angehören.

    \item Die Kassenprüfer prüfen die Buchführung, Jahresabrechnungen und die ordnungsgemäße Mittelverwendung des Vereins. Sie erstatten der Mitgliederversammlung Bericht und geben eine Empfehlung zur Entlastung des Vorstands.
\end{enumerate}

\refstepcounter{paragraphcounter}
\subsection{§ \theparagraphcounter\ Auflösung des Vereins}
\begin{enumerate}
    \item Der Verein löst sich auf, wenn die Mitgliederzahl unter sieben herabsinkt oder nach Beratung in zwei aufeinander folgenden, außerordentlichen Hauptversammlungen, in deren Tagesordnung die Beschlussfassung über eine Auflösung des Vereins angekündigt wurde. Für eine Auflösung müssen vier Fünftel der anwesenden, stimmberechtigten Mitglieder stimmen.

    \item Die Auflösung ist dem Amtsgericht in Stuttgart zur Löschung im Vereinsregister anzuzeigen.

    \item Im Falle einer Auflösung bzw. Aufhebung des Vereins, geht das Vereinsvermögen, nach Abzug der bestehenden Verbindlichkeiten, in die Verwaltung der Stadt Ostfildern über. Die Stadt hat das Vermögen so lange zu verwalten, bis sich in Ostfildern wieder ein Narrenverein mit der in § 2 genannten Zweckbestimmung bildet. Ehemalige Mitglieder des Narrenvereins - Narrenbund Ostfildern e. V. - müssen ein Recht auf Aufnahme in diesen Verein haben. Danach hat die Stadt das verwaltete Vermögen Letzterem auszuhändigen. Sollte nach Ablauf von 10 Jahren die o. g. Voraussetzung nicht vorliegen, kann die Stadt Ostfildern das verwaltete Vereinsvermögen an das DRK Ostfildern, Ortsgruppe Nellingen zuführen.

    \item Beschlüsse, die die Verwendung des Vereinsvermögens bei Auflösung bzw. bei Aufhebung des Vereins oder bei Wegfall seines bisherigen Zwecks beinhalten, dürfen erst nach Einwilligung des Finanzamtes ausgeführt werden.

    \item Sofern die Mitgliederversammlung nichts anders beschließt, sind der 1. Vorstand und der 1. stellv. Vorstand gemeinsam Vertretungsberechtigte Liquidatoren.
\end{enumerate}

\refstepcounter{paragraphcounter}
\subsection{§ \theparagraphcounter\ Ordnungen}
\begin{enumerate}
    \item Der Verein kann zur Regelung interner Abläufe Ordnungen erlassen. Diese Ordnungen dürfen der Satzung nicht widersprechen.

    \item Der Vereinsausschuss erarbeitet und beschließt eine Geschäftsordnung. Die Geschäftsordnung bedarf zur Wirksamkeit der Genehmigung durch die Mitgliederversammlung. Änderungen und die Aufhebung der Geschäftsordnung erfolgen im gleichen Verfahren.

    \item Der Vereinsausschuss ist ermächtigt, weitere Ordnungen zu erlassen, zu ändern oder aufzuheben, soweit die Satzung keine andere Zuständigkeit festlegt.

    \item Die Mitgliederversammlung kann eine Beitragsordnung beschließen, in der Beitragsklassen und Zahlungsmodalitäten geregelt werden.

    \item Von der Mitgliederversammlung beschlossene Ordnungen dürfen nur durch die Mitgliederversammlung geändert oder aufgehoben werden.
\end{enumerate}

\refstepcounter{paragraphcounter}
\subsection{§ \theparagraphcounter\ Die Gruppen innerhalb der Zunft}
\begin{enumerate}
    \item Dem Narrenbund Ostfildern e.V. gehören verschiedene Gruppen an.

    \item Sämtliche Gruppen unterstehen der Vorstandschaft des Vereins.

    \item Die Besetzung der Ausschussämter soll nach Möglichkeit so erfolgen, dass die im Verein bestehenden Gruppen im Ausschuss vertreten sind. Ein Anspruch der Gruppen auf Entsendung oder ein festes Mandat besteht jedoch nicht.

    \item Die Untergruppen können interne Ausschüsse bilden, die dem Zweck der kameradschaftlichen Einheit dienen, die Probleme jeder Art zu lösen versuchen und auch mit Aufgaben vertraut gemacht werden können, die den Vorstand entlasten.

    \item Beschlüsse und Satzung können diese Ausschüsse nicht anstreben und auch nicht geben. Diese Ausschüsse und ihre Personen haben im Vorstand kein Stimmrecht.
\end{enumerate}

\refstepcounter{paragraphcounter}
\subsection{§ \theparagraphcounter\ Inkrafttreten}
\begin{enumerate}
    \item Diese Satzung vom dd. MMMM YYYY – beschlossen in der Mitgliederversammlung vom dd. MMMM YYYY - ersetzt alle vorausgegangenen Satzungen und wird mit Eintragung ins Vereinsregister wirksam.
\end{enumerate}